% Generated draft: requires mathematical and source review.
\documentclass[11pt]{article}
\usepackage[margin=25mm]{geometry}
\usepackage{fontspec}
\setmainfont{Latin Modern Roman}
\usepackage{amsmath,amssymb,mathtools,unicode-math}
\setmathfont{Latin Modern Math}
\usepackage{enumitem,needspace,xurl,hyperref}
\setlength{\parindent}{0pt}
\setlength{\parskip}{5pt}
\setlength{\emergencystretch}{3em}
\newcommand{\N}{\mathbb{N}}
\newcommand{\Z}{\mathbb{Z}}
\newcommand{\Q}{\mathbb{Q}}
\newcommand{\R}{\mathbb{R}}
\newcommand{\C}{\mathbb{C}}
\newcommand{\F}{\mathbb{F}}
\newcommand{\A}{\mathbb{A}}
\newcommand{\PP}{\mathbb{P}}
\newcommand{\E}{\mathbb{E}}
\newcommand{\eps}{\varepsilon}
\newcommand{\dd}{\,\mathrm{d}}
\begin{document}
\section*{UnsolvedMath ID: 701397}
\textbf{Source problem code:} catalogue entry
\subsection{Definitions and mathematical statement}
Let \(G\leq\operatorname{Sym}(\Omega)\), where \(\Omega\) is a finite set. A partition \(\rho\) of \(\Omega\) is a collection of pairwise disjoint nonempty blocks whose union is \(\Omega\). A transversal of \(\rho\) is a subset meeting each block in exactly one point. For \(T\subseteq\Omega\), write \(T^g=\{t^g:t\in T\}\). The partition \(\rho\) is \(G\)-regular if
\[
\exists T\subseteq\Omega\quad\forall g\in G\quad\forall B\in\rho,\qquad |T^g\cap B|=1.
\]
The two trivial partitions are the partition into singletons and the partition with just one block. Thus a non-trivial proper partition has \(1<|\rho|<|\Omega|\). The group \(G\) is synchronizing if \(|\Omega|>2\) and there is no \(G\)-regular partition satisfying these inequalities. Non-synchronizing means that this definition fails; when \(|\Omega|>2\), this is equivalent to the existence of such a partition.

\textbf{(a)} Determine whether the following groups are synchronizing in the primitive affine actions intended by the source:
\[
2^p.\mathrm{PSL}(2,2p+1),
\qquad p\text{ and }2p+1\text{ prime},\quad
p\equiv3\pmod4,\quad p>23,
\]
and
\[
2^{101}.He.
\]
Here the notation \(2^d\) denotes an elementary abelian group of order \(2^d\), identified with the additive group of \(V=\F_2^d\). To specify a primitive affine action with linear group \(H\), one needs a faithful irreducible representation \(\alpha:H\longrightarrow\operatorname{GL}(V)\); the associated affine permutation group consists of the maps
\[
x\longmapsto v+\alpha(h)x\qquad(v\in V,\ h\in H).
\]
The intended pairs \((d,H)\) are \((p,\mathrm{PSL}(2,2p+1))\), subject to the displayed restrictions, and \((101,He)\). The supplied record does not specify the representations or embeddings. Consequently, these are questions about the intended actions, whose precise identification remains unresolved; the record does not justify replacing them by questions quantified over every faithful irreducible representation. The notation \(2^{101}.He\), including its exponent, is retained exactly as supplied.

\textbf{(b)} For each finite simple group \(S\), let
\[
G_S=\{f_{g,h}:g,h\in S\}\leq\operatorname{Sym}(S),
\qquad f_{g,h}(x)=g^{-1}xh.
\]
Classify the finite simple groups \(S\) for which this action of \(S\times S\) is non-synchronizing. Equivalently, for \(|S|>2\), classify those \(S\) for which there exist a partition \(\rho\) of the underlying set of \(S\) and a subset \(T\subseteq S\) such that
\[
1<|\rho|<|S|,
\qquad
\forall g,h\in S\quad\forall B\in\rho,\qquad
\bigl|g^{-1}Th\cap B\bigr|=1,
\]
where \(g^{-1}Th=\{g^{-1}th:t\in T\}\). The source says finite simple groups without an explicit nonabelian restriction. Under that literal scope, the group of order \(2\) is non-synchronizing because the degree condition fails. The notation \(G_S\) denotes the induced permutation group, so no faithfulness assumption on the abstract \(S\times S\)-action is added.
\subsection{Short English statement}
P. M. Neumann asks whether two specified families of primitive affine permutation groups are synchronizing, and which finite simple groups yield a non-synchronizing left-right action on themselves. Synchronization excludes partitions admitting a subset whose every group translate is a transversal. The affine representations needed to identify the actions in part (a) are not supplied. This is a statement-only draft requiring human mathematical review.
\subsection{Sources}
\begin{itemize}
\item Supplied local UnsolvedMath registry record: permanent ID \(701397\), exact source problem code \texttt{catalogue entry}, title \emph{Kourovka Notebook Problem 17.95}; upstream revision \texttt{d371ee917ac9647fbb06cd6fc529deec6f3dda82}.
\item Attributed in the supplied record to P. M. Neumann, proposed in 2010. The stated source is \emph{The Kourovka Notebook (Unsolved Problems in Group Theory)}, 17th issue (2010), original problem number 17.95, with the 21st issue edition (2026) also identified.
\item Supplied edition identifier: \texttt{arXiv:1401.0300v48}, \url{https://arxiv.org/abs/1401.0300v48}. The registry describes its collection as covering Issues 1--21 (1965--2026), edited by E. I. Khukhro and V. D. Mazurov. These bibliographic details are supplied attribution and have not been independently verified.
\item The registry labels this problem \emph{open}. This is a supplied status label, not an independently verified assessment of the literature. No separate literature assessment, literature-verification date, or additional literature sources were supplied.
\end{itemize}
\end{document}
